%% Pôster A0 da 70ª RBras — GAMLSS para produção diária de leite.
%% Compile com: pdflatex poster.tex
%% Dados (centis.csv, pontos.csv, familias.csv) vêm de dados_poster.R,
%% rodado a partir da raiz do projeto.

\documentclass{rbras-poster}

% Título e autores longos: corpo reduzido só aqui, para caber no cabeçalho.
\titulo{\fontsize{60bp}{66bp}\selectfont Beyond the average lactation curve:
  distributional regression for daily milk yield in Holstein cows}
\autores{\fontsize{28bp}{34bp}\selectfont \textbf{João Paulo Assis Bonifácio} · Pedro Mambelli Fernandes\\
  Geovane Fernandes da Silva · Gean Pereira Damaceno\\
  João Vitor Andrade Alves de Souza · Nayara Maria Barbosa de Sousa\\
  Marina de Arruda Camargo Danes · Luiz Ricardo Nakamura}
\afiliacoes{Universidade Federal de Lavras (UFLA), Lavras, MG}
% E-mails na ordem dos autores, em corpo menor para caber no cabeçalho.
\email{\fontsize{22bp}{27bp}\selectfont
  jpab.27@hotmail.com · pedromambelli@gmail.com · geovanefs90@gmail.com\\
  geandamaceno030@gmail.com · joaovitor.alvez16@gmail.com ·
  nayara.sousa@estudante.ufla.br\\ marina.danes@ufla.br · luiz.nakamura@ufla.com}
\mensagem{No pico da lactação, a vaca mediana produz 31,5~L/dia, mas 90\%
  do rebanho fica entre 17,8 e 46,8~L.}
% QR code do repositório, com correção de erro alta (H, 30%) para o logo central.
\qrset{level=H}
\qrlink{https://github.com/bonijoao/rbras70-curvas-centilicas-leite}
\qrlegenda{Código e pôster no GitHub}
% Logo do GitHub no centro do QR (centro do quadro: 2093, 790 bp do topo).
\AddToShipoutPictureFG*{\AtPageUpperLeft{%
  \begin{tikzpicture}[x=1bp,y=-1bp,overlay]
  \fill[white] (2093,790) circle (40bp);
  \node[inner sep=0] at (2093,790)
    {\includegraphics[width=64bp]{github-mark.pdf}};
  \end{tikzpicture}}}
\logoevento{logo-rbras.png}
\logoinstituicao{logo-ufla.png}
\conclusoes
{Os centis mostram se uma vaca está atípica para o seu DEL.}
{Dispersão, assimetria e curtose mudam ao longo da lactação.}
{Próximo passo: efeito de vaca e vários rebanhos.}
\referencias{[1] Rigby, R. A.; Stasinopoulos, D. M. (2005). Generalized additive
  models for location, scale and shape. \textit{JRSS C}, 54, 507--554.\par
  [2] Rigby, R. A.; Stasinopoulos, D. M. (2004). Smooth centile curves for skew
  and kurtotic data modelled using the Box-Cox power exponential distribution.
  \textit{Statistics in Medicine}, 23, 3053--3076.}
\agradecimentos{Fazenda Palmital (UFLA), pela cessão dos dados.}

\begin{document}
\hyphenpenalty=10000 \exhyphenpenalty=10000

\begin{colunaesquerda}
\begin{secao}{1}{Pergunta}
A curva média de lactação não diz se uma vaca produz pouco
\textit{para o seu estágio}.
\par\vspace{16bp}
\textbf{Como é a distribuição inteira da produção em cada dia de lactação?}
\end{secao}

\begin{secao}{2}{Dados}
\begin{itemize}
\item 62 vacas Holandesas, ordenha robótica
\item Fazenda Palmital (UFLA), Ijaci, MG
\item 9.514 registros diários, abr/2024 a jan/2025
\item Dias em lactação (DEL) de 1 a 705
\end{itemize}
\end{secao}

\begin{secao}{3}{Método}
GAMLSS: cada parâmetro da distribuição é uma curva suave do DEL.
\begin{caixaequacao}
\begin{equation*}
Y \mid \mathrm{DEL} \sim \mathrm{BCPE}(\mu,\sigma,\nu,\tau)
\end{equation*}
\begin{equation*}
g_k(\theta_k) = s_k(\mathrm{DEL}), \quad k = 1,\dots,4
\end{equation*}
{\tamRotulo\color{rbCinza} $\mu$ mediana · $\sigma$ dispersão\\
  $\nu$ assimetria · $\tau$ curtose\\ $s_k$: P-splines\par}
\end{caixaequacao}
\par\vspace{24bp}
Família escolhida pelo AIC; ajuste conferido por \textit{worm plot}.
\end{secao}
\end{colunaesquerda}

\begin{colunadireita}
\begin{secao}{4}{Resultados}
\duasfiguras{%
  \figura{O pico vem aos 60 DEL, com 29~L entre C5 e C95}
  {Produção diária de leite (L), por DEL}
  {\begin{axis}[rbras linha, xmin=0, xmax=400, ymin=0, ymax=60, clip mode=individual,
                 xtick={100,200,...,400}, ytick={0,10,...,60},
                 yticklabel style={/pgf/number format/.cd, fixed, precision=0}]
    \addplot[only marks, mark=*, mark size=1.6bp, rbGrade]
      table[x=DIM, y=y, col sep=comma] {pontos.csv};
    \foreach \c in {C5,C95}
      \addplot[rbCinzaClaro, line width=2.5bp]
        table[x=DIM, y=\c, col sep=comma] {centis.csv};
    \foreach \c in {C25,C75}
      \addplot[rbCinza, line width=3bp]
        table[x=DIM, y=\c, col sep=comma] {centis.csv};
    \addplot[rbAzul, line width=5.6bp]
      table[x=DIM, y=C50, col sep=comma] {centis.csv};
    \addplot[only marks, mark=*, mark size=7bp, rbAzul] coordinates {(60,31.46)};
    \node[anchor=south, font=\tamEixo\bfseries, text=rbAzul]
      at (axis cs:60,51) {31,5};
    \draw[rbAzul, line width=2bp] (axis cs:60,50.5) -- (axis cs:60,33.5);
    \node[anchor=west, font=\tamEixo, text=rbCinza] at (axis cs:405,32.2) {C95};
    \node[anchor=west, font=\tamEixo, text=rbCinza] at (axis cs:405,24.4) {C75};
    \node[anchor=west, font=\tamEixo\bfseries, text=rbAzul] at (axis cs:405,19.6) {C50};
    \node[anchor=west, font=\tamEixo, text=rbCinza] at (axis cs:405,14.6) {C25};
    \node[anchor=west, font=\tamEixo, text=rbCinza] at (axis cs:405,7.7) {C5};
    \end{axis}}
  {Centis 5, 25, 50, 75 e 95 do modelo BCPE. Pontos: amostra de 2.500 registros.}}
{%
  \figura{A faixa de 90\% encolhe de 32~L para 24~L}
  {Amplitude entre centis (L), por DEL}
  {\begin{axis}[rbras linha, xmin=0, xmax=400, ymin=0, ymax=35,
                 xtick={100,200,...,400}, ytick={0,5,...,35},
                 yticklabel style={/pgf/number format/.cd, fixed, precision=0}]
    \addplot[rbCinzaClaro, line width=3.75bp]
      table[x=DIM, y expr=\thisrow{C75}-\thisrow{C25}, col sep=comma] {centis.csv};
    \addplot[rbAzul, line width=5.6bp]
      table[x=DIM, y expr=\thisrow{C95}-\thisrow{C5}, col sep=comma] {centis.csv};
    \node[anchor=south, font=\tamEixo\bfseries, text=rbAzul]
      at (axis cs:230,28) {C95 $-$ C5};
    \node[anchor=south, font=\tamEixo, text=rbCinza]
      at (axis cs:300,9.8) {C75 $-$ C25};
    \end{axis}}
  {Variância constante, como nos modelos clássicos, não descreveria o rebanho.}}

\vspace{90bp}

\tabela{Ajustar a curtose melhora o modelo: a BCPE vence}
{\renewcommand\arraystretch{1.5}%
  \setlength\aboverulesep{0pt}\setlength\belowrulesep{0pt}%
  \begin{tabularx}{\linewidth}{@{}>{\raggedright\arraybackslash}X
                                 *{3}{>{\raggedleft\arraybackslash}X}@{}}
  \textbf{Distribuição} & \textbf{Parâmetros} & \textbf{AIC} & \textbf{$\Delta$AIC}\\
  \arrayrulecolor{rbAzul}\specialrule{3bp}{0pt}{0pt}
  \rowcolor{rbFundo}
  \textbf{BCPE} & \textbf{4} & \textbf{64.301} & \textbf{0}\\
  \arrayrulecolor{rbGrade}\specialrule{1bp}{0pt}{0pt}
  Box-Cox \textit{t} & 4 & 64.351 & 49\\
  \specialrule{1bp}{0pt}{0pt}
  Box-Cox Cole-Green & 3 & 64.441 & 140\\
  \specialrule{1bp}{0pt}{0pt}
  \end{tabularx}}
{Todos os parâmetros como P-splines do DEL.}
\end{secao}
\end{colunadireita}

\end{document}
